\documentclass[12pt,a4paper,openany,oneside]{book}

% --- Paquetes ---
\usepackage[utf8]{inputenc}
\usepackage[spanish, es-tabla]{babel}
\usepackage{amsmath, amsfonts, amssymb}
\usepackage{graphicx}
\usepackage{geometry}
\usepackage{hyperref}
\usepackage{booktabs}
\usepackage{fancyhdr}
\usepackage{listings}
\usepackage{xcolor}
\usepackage{tikz}
\usetikzlibrary{arrows.meta, positioning, shapes.geometric}

\geometry{margin=2.5cm}

% --- Estilo de Código Python ---
\definecolor{codegreen}{rgb}{0,0.6,0}
\definecolor{codegray}{rgb}{0.5,0.5,0.5}
\definecolor{codepurple}{rgb}{0.58,0,0.82}
\definecolor{backcolour}{rgb}{0.95,0.95,0.92}

\lstset{
    language=Python,
    backgroundcolor=\color{backcolour},   
    commentstyle=\color{codegreen},
    keywordstyle=\color{blue},
    numberstyle=\tiny\color{codegray},
    stringstyle=\color{codepurple},
    basicstyle=\ttfamily\small,
    breaklines=true,
    frame=single,
    showstringspaces=false
}

% --- Cabecera y Pie de Página ---
\pagestyle{fancy}
\fancyhf{}
\renewcommand{\headrulewidth}{0.4pt}
\renewcommand{\footrulewidth}{0.4pt}
\fancyhead[L]{\small Carlos César Sánchez Coronel}
\fancyhead[R]{\small Sesión 1: SGBD e Infraestructura de Datos}
\fancyfoot[L]{\small Carlos César Sánchez Coronel}
\fancyfoot[C]{\thepage}

% --- Portada ---
\title{
    \textbf{Sesión 0} \\ 
    \Huge \textbf{SYLLABUS A BBDD }
}
\author{
    \small Por \\ \vspace{0.2cm}
    \Large \textsc{Carlos César Sánchez Coronel}
}
\date{2026}

\begin{document}

\maketitle
\tableofcontents






\chapter{Fundamentos, Diseño y Gobernanza de Datos (Semanas 1-4)}

\section{Semana 1: El Universo del Dato y la Infraestructura Física}

\subsection{La Materia Prima: Del Dato a la Metadata}
Para un ingeniero, la gestión de datos implica controlar la entropía de la información. Debemos distinguir claramente:
\begin{itemize}
    \item \textbf{Dato:} Valor atómico crudo (ej. $1010$).
    \item \textbf{Información:} Dato con significado (ej. ID de cliente $1010$).
    \item \textbf{Metadata:} Datos que describen otros datos. En IA es crítico para el \textit{Data Lineage} (linaje) y asegurar que el set es apto para entrenamiento.
    \item \textbf{Logs y Temporales:} Registros de eventos para auditoría y archivos de caché para optimizar la velocidad de respuesta.
\end{itemize}

\subsection{Hardware: El Soporte Físico y Procesos de Búsqueda}
El rendimiento de una BBDD depende de la interacción entre el software y el hardware:
\begin{itemize}
    \item \textbf{Comportamiento en Hardware:}
    \begin{itemize}
        \item \textbf{Buffer Pool (RAM):} El motor busca primero en RAM. Si hay un \textit{Cache Hit}, el acceso es $O(1)$ en nanosegundos.
        \item \textbf{I/O de Disco (SSD/NVMe):} Indispensables para producción. Si el dato no está en RAM, el hardware sufre latencia por el bus de datos hacia el disco.
        \item \textbf{SAN (Storage Area Network):} Redes dedicadas para almacenamiento masivo y redundancia.
    \end{itemize}
    \item \textbf{Complejidad Computacional de Búsqueda:}
    \begin{itemize}
        \item \textbf{Table Scan:} Lectura secuencial, complejidad $O(n)$. Inviable en grandes volúmenes.
        \item \textbf{Índices B-Tree:} Estándar en SQL, complejidad $O(\log n)$.
        \item \textbf{Hash Indexing:} Para búsquedas de igualdad exacta, logrando $O(1)$.
    \end{itemize}
\end{itemize}



\section{Semana 2: Roles y Marcos de Gobernanza}

\subsection{Roles en la Gestión de Datos}
\begin{itemize}
    \item \textbf{DBA (Database Administrator):} Responsable de salud del hardware, instalación, parches de seguridad (especialmente en Oracle/SQL Server) y tuning.
    \item \textbf{Analista / Data Engineer:} Diseña modelos lógicos, optimiza el costo computacional de las consultas y construye pipelines para IA usando Python.
\end{itemize}

\subsection{Gobernanza y Explotación (Data Governance)}
Es el marco normativo que asegura que los datos sean confiables, privados y seguros:
\begin{itemize}
    \item \textbf{Calidad de Datos:} Reglas para asegurar que el dato sea preciso y completo (\textit{Garbage In, Garbage Out}).
    \item \textbf{Privacidad y Seguridad:} Cumplimiento de normativas (GDPR/Protección de Datos Personales) y enmascaramiento de datos sensibles.
    \item \textbf{Linaje de Datos (Data Lineage):} Capacidad de rastrear el origen de un dato y sus transformaciones hasta que llegó al modelo de IA.
\end{itemize}


\subsection{Panorama de Motores Líderes en la Industria}
\begin{itemize}
    \item \textbf{Oracle Database:} El referente histórico en el mercado corporativo, destacando por su robusta gestión de memoria y alta disponibilidad en transacciones críticas.
    \item \textbf{PostgreSQL:} La opción predilecta para proyectos de Inteligencia Artificial debido a su flexibilidad y soporte de extensiones como \texttt{pgvector}.
    \item \textbf{SQL Server y MySQL:} Motores de alta fiabilidad líderes en la integración con ecosistemas empresariales de Microsoft y aplicaciones web respectivamente.
\end{itemize}

\subsection{Conectividad y Visualización (BI)}
Los SGBD se conectan a herramientas de Business Intelligence mediante conectores estándar (ODBC/JDBC):
\begin{itemize}
    \item \textbf{Power BI y Tableau:} Conexión directa a SQL Server, Oracle o BigQuery para dashboards dinámicos.
    \item \textbf{Streamlit / Dash:} Librerías de Python para crear interfaces rápidas que consumen datos del SGBD y muestran resultados de modelos de IA.
\end{itemize}

\section{Semana 3: Modelamiento y Arquitectura Conceptual}

\subsection{Ciclo de Modelamiento: Del Concepto al Disco}
El diseño requiere niveles de abstracción para garantizar que la estructura soporte las reglas de negocio:
\begin{itemize}
    \item \textbf{Modelo Entidad-Relación (E-R):} Nivel conceptual. Se identifican entidades (sustantivos), atributos (características) y relaciones (verbos). Es independiente del software.
    \item \textbf{Modelo Lógico:} Traduce el E-R a estructuras de tablas, definiendo Llaves Primarias (PK) y Llaves Foráneas (FK). Se aplica la normalización.
    \item \textbf{Modelo Físico:} Implementación real en el SGBD. Incluye tipos de datos (VARCHAR, INT), índices, particionamiento y ubicación en el disco.
\end{itemize}



\section{Semana 4: Normalización e Integridad Técnica}

\subsection{Normalización: Control de la Redundancia}
Aplicamos las Formas Normales (FN) para evitar anomalías de inserción y borrado:
\begin{itemize}
    \item \textbf{1FN (Atomicidad):} Valores atómicos; no se permiten grupos repetidos ni listas en una celda.
    \item \textbf{2FN (Dependencia Total):} En 1FN y cada columna no llave depende de toda la llave primaria.
    \item \textbf{3FN (Dependencia Transitiva):} En 2FN y las columnas no llave no deben depender de otras columnas no llave.
\end{itemize}



\subsection{Integridad Referencial y Restricciones}
\begin{itemize}
    \item \textbf{Primary Key (PK):} Identificador único e irreducible de una fila.
    \item \textbf{Foreign Key (FK):} Vincula una tabla con la PK de otra. Garantiza que no existan "hijos huérfanos".
    \item \textbf{Constraints (Check/Unique):} Reglas de validación a nivel de motor (ej. sueldo $> 0$).
    \item \textbf{Dimensionamiento Lógico:} Un mal diseño obliga al motor a realizar \texttt{JOINS} excesivos, elevando el uso de CPU y RAM, crítico en volúmenes de \textbf{PB} en banca.
\end{itemize}








\chapter{Dominio del Lenguaje SQL y Lógica de Servidor (Semanas 5-7)}

\section{Semana 5: Manipulación de Datos y Lógica de Conjuntos}

\subsection{El Motor de Operatividad: DML}
El lenguaje de manipulación de datos (DML) permite interactuar con el estado de los registros y es el núcleo de cualquier analista o desarrollador:
\begin{itemize}
    \item \textbf{SELECT:} Proyección de columnas y filtrado de filas mediante predicados lógicos.
    \item \textbf{INSERT / UPDATE / DELETE:} Operaciones transaccionales para la creación, modificación y eliminación de la persistencia de datos.
\end{itemize}

\subsection{Relaciones mediante JOINS}
Para un ingeniero, los \textit{Joins} representan operaciones lógicas de conjuntos necesarias para reconstruir información que ha sido previamente normalizada:
\begin{itemize}
    \item \textbf{INNER JOIN:} Obtiene la intersección de conjuntos donde existe coincidencia en ambos universos ($A \cap B$).
    \item \textbf{LEFT / RIGHT JOIN:} Permite preservar la integridad de un universo (tabla maestra) frente a sus relaciones opcionales en tablas secundarias.
    \item \textbf{FULL OUTER JOIN:} Realiza la unión completa de los universos de datos involucrados ($A \cup B$).
\end{itemize}



\section{Semana 6: Análisis Avanzado y Transformación de Datos}

\subsection{Agregaciones y Agrupamientos}
Uso de la cláusula \texttt{GROUP BY} para resumir grandes volúmenes de información. Esto permite aplicar funciones estadísticas esenciales como \texttt{SUM}, \texttt{AVG}, \texttt{COUNT}, \texttt{MIN} y \texttt{MAX}, junto con el filtrado avanzado post-agregación mediante \texttt{HAVING}.

\subsection{Funciones de Ventana (Window Functions)}
En el contexto de IA y análisis de series temporales, las funciones de ventana permiten realizar cálculos complejos (ej. \texttt{OVER}, \texttt{RANK}, \texttt{ROW\_NUMBER}, \texttt{LEAD}, \texttt{LAG}) sobre un conjunto de filas relacionadas sin necesidad de colapsar el resultado en un solo grupo, manteniendo el detalle de cada fila.

\subsection{Limpieza de Datos (Data Cleaning)}
Dominio de funciones especializadas para garantizar la calidad del dato antes de su procesamiento:
\begin{itemize}
    \item \textbf{Texto:} Operaciones de transformación con \texttt{UPPER}, \texttt{SUBSTR}, \texttt{CONCAT}, \texttt{TRIM} y el uso de \texttt{LIKE} o expresiones regulares (\textit{RegEx}).
    \item \textbf{Fecha/Hora:} Gestión de zonas horarias, aritmética de intervalos y extracción de componentes temporales (año, mes, día).
\end{itemize}

\section{Semana 7: Objetos de Lógica Interna y Mercado de Motores}

\subsection{Programación en el Servidor}
Los SGBD modernos permiten programar lógica de negocio cerca de los datos, lo que minimiza drásticamente la latencia de red:
\begin{itemize}
    \item \textbf{Vistas (Views):} Tablas virtuales que encapsulan consultas complejas, simplificando el acceso para el usuario final y añadiendo capas de seguridad.
    \item \textbf{Funciones y Procedimientos Almacenados:} Bloques de código reutilizables que permiten automatizar procesos y cálculos pesados directamente en el motor.
    \item \textbf{Triggers (Disparadores):} Automatismos que se ejecutan ante eventos DML. Son fundamentales en el sector bancario para generar \textbf{Logs de Auditoría} y mantener la consistencia de saldos en tiempo real.
\end{itemize}
















\chapter{Ingeniería de Datos y Big Data (Semanas 8-10)}

\section{Semana 8: Captura de Datos y Extracción Masiva}

\subsection{Fuentes de Datos y Web Scraping}
Para alimentar modelos de IA modernos, los datos deben ser extraídos de diversas fuentes externas e internas:
\begin{itemize}
    \item \textbf{Web Scraping:} Técnica fundamental para extraer datos directamente de sitios web. Fue la fuente principal para entrenar LLMs modernos usando datasets como \textit{Common Crawl}.
    \item \textbf{Wayback Machine (Internet Archive):} Recurso crítico para recuperar datos históricos de la web que ya no están disponibles en línea.
\end{itemize}

\subsection{Captura de Datos en Tiempo Real}
\begin{itemize}
    \item \textbf{APIs:} Interfases de programación para el intercambio de datos estructurados entre sistemas.
    \item \textbf{CDC (Change Data Capture):} Métodos para extraer cambios en sistemas transaccionales en tiempo real sin afectar el rendimiento del SGBD principal.
\end{itemize}



\section{Semana 9: Pipelines y Paradigmas de Procesamiento}

\subsection{ELT vs. ETL}
La evolución del almacenamiento en la nube ha cambiado la forma en que movemos los datos:
\begin{itemize}
    \item \textbf{ETL (Extract, Transform, Load):} Los datos se transforman en un servidor intermedio antes de llegar al destino. Es ideal para datos estructurados y limpieza previa.
    \item \textbf{ELT (Extract, Load, Transform):} Los datos crudos se cargan directamente en el motor y se transforman usando su propio poder de cómputo. Es el estándar actual en Cloud e IA.
\end{itemize}

\subsection{Batch vs. Streaming}
\begin{itemize}
    \item \textbf{Batch (Lotes):} Procesamiento de grandes volúmenes en intervalos de tiempo, como el cierre diario en banca.
    \item \textbf{Streaming (Tiempo Real):} Procesamiento de eventos conforme ocurren, vital para la detección de fraude en segundos. Herramientas clave: Apache Kafka y AWS Kinesis.
\end{itemize}



\section{Semana 10: Big Data y Formatos de Alto Rendimiento}

\subsection{Almacenamiento Optimizado}
En entornos de IA, el almacenamiento no siempre es una tabla viva, sino archivos optimizados:
\begin{itemize}
    \item \textbf{Apache Parquet:} Formato columnar que permite lecturas rápidas de subconjuntos de datos. Es el formato nativo para procesos de entrenamiento masivo de modelos.
    \item \textbf{Data Lakes vs. Data Warehouses:} El primero guarda datos crudos (archivos Parquet/JSON) y el segundo datos estructurados.
\end{itemize}

\subsection{Plataformas de Big Data y Lakehouse}
\begin{itemize}
    \item \textbf{Big Data:} El reto de las 3 "V" (Volumen, Velocidad y Variedad). Se requiere herramientas como Apache Spark para procesar volúmenes que el SQL tradicional no puede manejar.
    \item \textbf{Databricks:} Plataforma que utiliza la arquitectura \textbf{Lakehouse}, uniendo la flexibilidad del Data Lake con el rigor del SGBD.
\end{itemize}








\chapter{Ecosistemas Modernos e IA Generativa (Semanas 11-12)}

\section{Semana 11: El Paradigma NoSQL y el Ecosistema Cloud}

\subsection{SGBD NoSQL y su Integración con IA}
Sistemas diseñados para alta escalabilidad y flexibilidad, integrables con Python mediante librerías como \texttt{pymongo}, \texttt{redis-py} y \texttt{cassandra-driver}:
\begin{itemize}
    \item \textbf{Documentales (MongoDB):} Almacenan documentos en formato JSON/BSON. Ideales para catálogos y datos con esquemas cambiantes.
    \item \textbf{Clave-Valor (Redis):} Bases de datos \textit{In-Memory} que ofrecen latencias mínimas. Cruciales para la gestión de estados y memoria de corto plazo en aplicaciones de IA.
    \item \textbf{Grafos (Neo4j):} Esenciales para modelar relaciones complejas, sistemas de recomendación y detección de fraude.
    \item \textbf{Columnares (Cassandra):} Diseñadas para manejar Petabytes de datos distribuidos con alta disponibilidad.
    \item \textbf{Motores de Búsqueda (Elasticsearch):} Especializados en la indexación y búsqueda de texto masivo.
\end{itemize}

\subsection{Infraestructura en la Nube: On-Premise vs. Cloud}
\begin{itemize}
    \item \textbf{On-Premise:} Control total del hardware y seguridad física, pero con una alta inversión inicial (CapEx).
    \item \textbf{Cloud Native (PaaS):} Servicios gestionados que permiten escalado elástico. Ejemplos líderes: \textbf{AWS} (Aurora, DynamoDB), \textbf{Azure} (SQL Database, Cosmos DB) y \textbf{GCP} (Cloud SQL, BigQuery).
\end{itemize}



\section{Semana 12: Bases de Datos Vectoriales y RAG}

Las bases de datos vectoriales son el pilar de la IA Generativa. A diferencia de las tradicionales, estas operan sobre la \textbf{semántica} y no solo sobre coincidencias exactas.

\subsection{Conceptos Fundamentales de Vectores}
\begin{itemize}
    \item \textbf{Embeddings:} Representaciones numéricas de datos (texto, imágenes) en un espacio de alta dimensionalidad. Un modelo de IA traduce una frase en un vector donde la posición relativa indica su significado.
    \item \textbf{Búsqueda por Similitud (ANN):} Algoritmos de \textit{Approximate Nearest Neighbors} (como HNSW o IVF) que encuentran los vectores más cercanos en milisegundos.
    \item \textbf{Métricas de Distancia:} Cálculo de la "cercanía" semántica mediante funciones como la Similitud de Coseno o la Distancia Euclidiana.
\end{itemize}

\subsection{Flujo de Trabajo: Retrieval Augmented Generation (RAG)}
El patrón RAG soluciona las alucinaciones de los LLMs conectándolos a datos reales:
\begin{enumerate}
    \item El usuario realiza una consulta.
    \item La consulta se convierte en un vector (\textit{Query Embedding}).
    \item La BBDD Vectorial recupera los fragmentos de información más parecidos.
    \item El LLM genera una respuesta utilizando esos fragmentos como contexto verídico.
\end{enumerate}



\subsection{Tecnologías y Orquestación}
\begin{itemize}
    \item \textbf{Nativas y Extensiones:} Herramientas como \textbf{Pinecone} o \textbf{Milvus} (nativas) y extensiones como \textbf{pgvector} para PostgreSQL.
    \item \textbf{Orquestadores:} Uso de \textbf{LangChain} y \textbf{LlamaIndex} para conectar la base de datos con los modelos de lenguaje en Python.
\end{itemize}




\chapter{Arquitectura, Resiliencia y Gestión de Desastres (Semanas 13-14)}

\section{Semana 13: Despliegue y Alta Disponibilidad}

\subsection{Ecosistema de Contenedores: SGBD en Docker}
En la ingeniería moderna, las BBDD se despliegan mediante contenedores para garantizar la consistencia entre entornos:
\begin{itemize}
    \item \textbf{Dockerización:} Empaquetamiento del SGBD en imágenes. Es imperativo el uso de \textbf{Volúmenes} para que la información persista independientemente del ciclo de vida del contenedor.
    \item \textbf{Orquestación con Kubernetes:} Gestión del auto-escalado y la disponibilidad del motor ante incrementos en la carga de peticiones.
\end{itemize}

\subsection{Estrategias de Disponibilidad (HA)}
Para evitar caídas de servicio, especialmente en sectores críticos, se implementan:
\begin{itemize}
    \item \textbf{Réplicas de Lectura:} Copias en tiempo real que permiten distribuir la carga. La IA consume datos de estas réplicas para no degradar el performance del nodo transaccional.
    \item \textbf{Failover:} Capacidad del sistema para conmutar automáticamente a un nodo secundario si el principal falla.
\end{itemize}



\section{Semana 14: Gestión de Errores y Recuperación (DRP)}

\subsection{Operaciones de Alto Riesgo}
El conocimiento de los comandos que alteran la estructura física es vital para evitar catástrofes:
\begin{itemize}
    \item \textbf{DROP vs. TRUNCATE:} El primero elimina la tabla y su metadata; el segundo vacía los datos de forma física sin generar logs detallados por fila.
    \item \textbf{UPDATE/DELETE sin WHERE:} El error humano más común. Se deben establecer restricciones y \textit{scripts} de validación previos en entornos de producción.
\end{itemize}

\subsection{Mecanismos de Resiliencia}
\begin{itemize}
    \item \textbf{Rollback:} Capacidad transaccional para revertir cambios ante una ejecución fallida, devolviendo la base de datos a un estado consistente.
    \item \textbf{Backups y Snapshots:} Diferenciación entre copias de seguridad lógicas (programadas) y fotos instantáneas del disco (Snapshots) para recuperación rápida.
    \item \textbf{Point-in-Time Recovery:} Uso de los \textit{Transaction Logs} para restaurar la base de datos a un segundo exacto previo a un incidente.
\end{itemize}



\subsection{Desafíos en la Nube y Data Drift}
\begin{itemize}
    \item \textbf{Latencia y Egreso:} Consideraciones sobre el costo y tiempo de mover Petabytes de datos entre regiones o hacia entornos locales.
    \item \textbf{Data Drift (Deriva de Datos):} En entornos de IA, un "error" también puede ser la degradación del modelo debido a cambios en la distribución de los datos, lo que requiere una gobernanza activa.
\end{itemize}

\section{Despliegue de Arquitecturas Híbridas (SQL, NoSQL y Vectorial)}

En proyectos de IA a gran escala, la arquitectura no se limita a un solo motor; se despliega un ecosistema de persistencia políglota donde cada base de datos cumple un rol específico según la naturaleza del dato y la necesidad de procesamiento.

\subsection{Orquestación de Persistencia Políglota}
El despliegue moderno utiliza contenedores (Docker) u orquestadores (Kubernetes) para levantar múltiples servicios que interactúan de forma coordinada:
\begin{itemize}
    \item \textbf{Capa Transaccional (SQL):} Motores como PostgreSQL u Oracle actúan como la fuente de verdad para los datos estructurados, maestros de clientes y registros operativos.
    \item \textbf{Capa de Velocidad (NoSQL):} Uso de Redis para el almacenamiento de sesiones de usuario, caché de inferencia y gestión de estados de corta duración.
    \item \textbf{Capa de Memoria Semántica (Vectorial):} Motores como Milvus, Pinecone o Weaviate para gestionar los embeddings, permitiendo realizar búsquedas por similitud y alimentar el flujo RAG.
\end{itemize}



\subsection{Estrategias de Despliegue Unificado}
\begin{itemize}
    \item \textbf{Infraestructura como Código (IaC):} Definición de redes, volúmenes y variables de entorno en archivos de configuración para garantizar que el despliegue sea replicable en cualquier nube o servidor local.
    \item \textbf{Aislamiento de Redes:} Configuración de redes privadas virtuales para que los motores de base de datos se comuniquen entre sí sin exposición directa a la red pública, aumentando la seguridad.
    \item \textbf{Gestión de Volúmenes Cruzados:} Implementación de almacenamiento persistente sincronizado para asegurar que todo el ecosistema mantenga la integridad de los datos ante un reinicio del sistema.
\end{itemize}



\subsection{Sincronización y Consistencia entre Motores}
El mayor reto técnico en este despliegue es mantener la paridad de información entre la base de datos relacional y la vectorial:
\begin{itemize}
    \item \textbf{Disparadores y CDC:} Uso de Triggers o herramientas de Change Data Capture para detectar inserciones en SQL y gatillar automáticamente la generación del Embedding y su inserción en la base de datos vectorial.
    \item \textbf{Consistencia Eventual:} Gestión de los tiempos de latencia desde que un dato se crea en el sistema transaccional hasta que está disponible para ser consultado semánticamente por el modelo de IA.
\end{itemize}


\end{document}